% =============================================================================
% SECTION 5: DISCUSSION
% =============================================================================
\section{Discussion}

\subsection{What the Framework Provides}

\textit{A standardized vocabulary.}
The notation $L_i.T_j$ gives precise loader descriptions.
For example, $(L_0.T_1, L_1.T_1, L_2.T_1, L_3.T_2, L_4.T_1, L_5.T_2)$ names a loader that uses anti-debug, \texttt{.rdata} storage, local RWX allocation, AES, \texttt{memcpy}, and \texttt{EnumDisplayMonitors} execution.
Comparing two loaders, or resolving terminology disagreements, reduces to comparing tuples.

\textit{A vehicle for controlled experiments.}
Because technique selection happens at preprocessor time, changing a single stage while keeping others constant produces a binary whose only code difference is inside that stage.
Detection differences between two such binaries can then be attributed to the changed stage with fewer confounding variables than comparing hand-written loaders.

\textit{Complementarity with MITRE ATT\&CK.}
ATT\&CK labels \emph{what} an attacker does at a strategic level; the pipeline model describes \emph{how} a specific loader works internally.
A loader classified under T1055 Process Injection in ATT\&CK can be further decomposed into its six stages and summarized as a $C$-tuple, which makes the internal mechanics explicit without displacing the ATT\&CK label.

\subsection{Limitations}

\begin{enumerate}[leftmargin=*, itemsep=2pt]
    \item \textit{Small technique inventory.}
    The platform currently provides 1--3 techniques per stage.
    Several categories that would be important for a broader survey are not yet implemented, including remote allocation, APC queuing, process hollowing, and network-based storage.

    \item \textit{Single security product.}
    Evaluation is limited to Windows Defender.
    Commercial EDR and other AV products use different telemetry and heuristics; results are not expected to transfer directly.

    \item \textit{Single, fixed shellcode.}
    We fix the shellcode to isolate loader effects on detection.
    This is a deliberate methodological choice, but it means the study does not address payload-level factors such as shellcode size, staged versus reflective payloads, or obfuscation inside the shellcode itself.
    These would require a separate axis of variation.

    \item \textit{Detection methodology is proposed, not validated.}
    The stage-aware detection methodology of Section~3.3 is a design.
    No rules have been written and no rule-triggering data exist.
    Its usefulness has to be established in subsequent work before claims can be made about coverage or gap detection.

    \item \textit{Non-stationary targets.}
    AV definitions and EDR heuristics change frequently, so a single-timepoint experiment captures a snapshot.
    Longitudinal studies are needed to assess stability.
\end{enumerate}

\subsection{Future Work}

\begin{itemize}[leftmargin=*, itemsep=2pt]
    \item \textit{Populate the detection methodology.}
    Write concrete Sigma and YARA rules for each technique in Table~\ref{tab:techniques}, populate the mapping $D$ from Eq.~(\ref{eq:detection_map}), and run the rules against telemetry already collected by the platform.
    This is the most direct step to turn Section~3.3 from a proposal into a validated method.

    \item \textit{Technique expansion.}
    Add remote injection (\texttt{VirtualAllocEx} $\rightarrow$ \texttt{WriteProcessMemory} $\rightarrow$ \texttt{CreateRemoteThread} or APC queuing), process hollowing, network-backed storage, and richer anti-analysis (timing checks, sandbox fingerprinting) so that the combinatorial space is large enough for meaningful comparison.

    \item \textit{Multi-product evaluation.}
    Integrate additional AV engines and at least one commercial EDR to build a cross-product detection matrix.
    The per-stage view is expected to be most useful here: the same loader configuration may produce different detection patterns across products, and the pipeline model can attribute those differences to specific stages.

    \item \textit{Real-world loader decomposition.}
    Apply the pipeline model to loaders extracted from public samples---Donut~\cite{donut}, ScareCrow~\cite{scarecrow}, Havoc~\cite{havoc} implants, and MalwareBazaar~\cite{malwarebazaar} submissions---by mapping each real loader onto a $C$-tuple.
    This tests the generality of the decomposition and builds a structured knowledge base of real-world configurations.

    \item \textit{Longitudinal tracking.}
    Re-run a fixed configuration set on a rolling schedule to track how detection outcomes shift as AV definitions and EDR models are updated.
\end{itemize}
